\textbf{Question 31.} \\
\hrule 

The composer Beethoven wrote 9 symphonies, 5 piano concertos, and 32 piano sonatas.

\textbf{a. How many ways are there to play first a Beethoven symphony and then a Beethoven piano concerto?} \newline 
We can just simply multiply the number of symphonies by the number of concertos in this case. So, $\boxed{9 \cdot 5 = 45}$ ways of organizing those compositions. 

\textbf{b. The manager of a radio station decides that on each successive evening (7 days per week), a Beethoven symphony will be played followed by a Beethoven piano concerto followed by a Beethoven piano sonata. For how many years could this policy be continued before exactly the same program would have to be repeated?} \newline 
Here we can do the same thing as part (b), but also multiply by the number of piano sonatas (32). Then, since we will have the number of days we can divide by 365 to roughly approximate the years this would last for. 
\[9 \cdot 5 \cdot 32 = \boxed{1440 \; \text{days} \; \approx \text{4 years}}\]
